\subsection{Efficiency, assumptions, and the path to a genuine ZK-KBFold PCS}\label{sec:zkrs:ledger}

\begin{definition}[Field-message accounting for the public core]\label{def:zkrs:count}
Count only publicly transmitted elements of $L$
in the core of Definition~\ref{def:zkrs:public-core},
not commitments, field traces, challenge bits or
the separate affine proofs. Let $n=|I|=2^\ell$.
\end{definition}

\begin{lemma}[Exact single-opening public-core message count]\label{lem:zkrs:count}
The honest public core transmits $2n+2$ elements of $L$:
$n$ entries of $m$, one $y'$, $n$ entries of $U$ and one $A$.
It also receives $\ell+1$ challenge elements of $L$:
the $\ell$ coordinates of $\theta$ and one $z$.
\end{lemma}
\begin{proof}
Read the five ordered stages of
Definition~\ref{def:zkrs:public-core}.
The lengths of the prover-message
vectors are $n,1,n,1$ respectively,
whose sum is $2n+2$.
The verifier supplies a vector
$\theta\in L^\ell$ and
a scalar $z\in L$, contributing
$\ell+1$ challenge elements.
No other elements are counted by
Definition~\ref{def:zkrs:count}.
\end{proof}

\begin{lemma}[Naive affine-coefficient count]\label{lem:zkrs:H-count}
The family of public affine coefficients
$H_{v,k}(\theta)$ from
Definition~\ref{def:binary:affine-H}
contains $n^2$ elements of $L$.
A direct implementation of their
defining sums uses $n$ summands
per coefficient, hence $n^3$
basis-coordinate terms.
\end{lemma}
\begin{proof}
Both indices $v,k$ range over
$I$, a set of size $n$, so
there are $n^2$ ordered pairs.
By definition, each coefficient
$H_{v,k}$ is a sum over a third
index $u\in I$ with $n$ possible
values. Counting the summands
gives $n\cdot n^2=n^3$ terms.
This is an operation-count upper
bound for the naive formula,
not a lower bound on optimized
implementations.
\end{proof}

\begin{lemma}[Exact private state field-element count]\label{lem:zkrs:hidden-count}
Excluding the already committed packed witness, the auxiliary
private tuple $(p,p_{\rm in},c,t,\rho,\sigma)$ has
$3n+3$ field slots: $n+1+n+n+1+1$.
In the honest sampler the genuinely independent
randomness in that tuple consists of
$2n+2$ field elements: $n$ for $p$,
one for $p_{\rm in}$, $n$ for $t$ and
one for $\rho$.
\end{lemma}
\begin{proof}
The displayed tuple has, in order,
$n,1,n,n,1,1$ field entries, summing to
$3n+3$. The honest sampler chooses
only $p,p_{\rm in},t,\rho$ independently.
Their entry counts sum to
$n+1+n+1=2n+2$.
The remaining entries satisfy
$c=\mathcal T(p)$ and
$\sigma=\rho^2$, so they
are deterministic functions
of the sampled entries.
\end{proof}

\begin{definition}[Certification milestones]\label{def:zkrs:milestones}
The following are \emph{separate}
certification tasks rather than
consequences of the preceding theorems:
\begin{enumerate}[label=\textup{(\roman*)}]
\item \textbf{Real commitment semantics:}
a concrete binding and hiding scheme
with a consistent simulator for
commitments and authenticated queries.
\item \textbf{Affine proof realization:}
an implementation that enforces every
relation in Definition~\ref{def:zkrs:affine-checks}
on one fixed hidden tuple with explicit
completeness, soundness and zero-knowledge error.
\item \textbf{Multi-opening ZK:}
commitment simulation before evaluation
points are known; fresh masks per opening;
and a joint simulator for every adaptive
challenge, proof transcript and opening.
\item \textbf{Cross-field integration:}
the source KBFold PCS uses odd-characteristic
multiplicative-subgroup word folding.
The present binary-field checker does not
inherit its proximity or QROM theorem.
\item \textbf{Concrete optimization:}
avoid or compress the naive $n^2$
public affine coefficients and verify
the resulting circuit/prover work,
rather than asserting asymptotic
efficiency without implementation.
\end{enumerate}
\end{definition}

\begin{proposition}[Binary-field folding is not directly interchangeable with KBFold's $\pm x$ fold]\label{prop:zkrs:cross-field}
The characteristic-two ring-switching checker
cannot reuse the original multiplicative
domain pairing $x,-x$ as a two-point
local fold over $L=\F_{2^n}$.
\end{proposition}
\begin{proof}
In characteristic two, $-x=x$
for every $x\in L$.
Thus the supposed two distinct
evaluation positions coincide.
Furthermore the standard
even/odd word-fold formulas
divide by $2x$, but $2=0$
in $L$, so this expression
has no inverse for any $x$.
The required field hypotheses
of Definition~\ref{def:kbff:wordfold}
therefore fail, as shown in
Lemma~\ref{lem:binary:no-even-subgroup}.
\end{proof}

\begin{remark}[Historical and originality boundary]
Diamond--Posen supply the original
tensor-algebra compiler
\cite{DiamondPosen2024}.
Kochergin's 2026 note
\cite{Kochergin2026}
and Frigo--shelat's
sumcheck bit-packing paper
\cite{FrigoShelat2026} give
different zero-knowledge routes
for binary-field packing.
The reusable mathematical contribution
developed here is the exact
kernel-coordinate/transposition
dictionary, its anchored affine
checker and a rigorous separation
between perfect public-core
simulation and full proof-system
simulation. This chapter does not
claim that trace-zero masking,
the general idea of replacing
quadratic witness constraints by
public squaring, or a
single-challenge polynomial identity
test originated here.
\end{remark}

\begin{remark}[The next theorem genuinely needed]
For the \emph{original} odd-characteristic
KBFold PCS, the central outstanding
theorem is a composable adaptive
view simulator for committed masked
fold oracles, including every
sumcheck round polynomial,
authenticated opening and the
degree-aware linear final closure
of Chapter~5.
A public-core theorem for binary
Ring Switching does not discharge
that security obligation.
\end{remark}
